\begin{Exercise}[label={ex:telefoni}]
    Una società di telefonia cellulare gestisce un programma di premiazione per "utenti fedeli". In particolare, per
    ogni cliente viene salvato su un file binario "punti.dat" il nome del cliente (al massimo 31 caratteri) e un
    numero intero che rappresenta i punti accumulati. Tali informazioni sono organizzate come una struttura \texttt{User}, opportunamente definita.

    Si scriva una funzione \texttt{readPoints()} che ricevuto in ingresso il nome di un file \texttt{usersFile}, un array \texttt{results} di
    strutture \texttt{User}, la dimensione massima dell’array \texttt{maxDim}, e un limite
    inferiore di punti \texttt{minPoints}, copi nell’array \texttt{results} i dati dei clienti che hanno almeno i punti specificati da \texttt{minPoints}.

    La funzione deve restituire come risultato il numero di utenti con almeno
    \texttt{minPoints}. Si noti che tale risultato rappresenta anche la dimensione
    logica dell’array results. Qualora il file non sia accessibile, la funzione deve restituire il valore -1.

    Si scriva poi un programma \texttt{main()} che chieda all’utente il numero
    di clienti salvati sul file (tale numero sarà noto solo a tempo di
    esecuzione), e allochi dinamicamente un vettore V di \texttt{User}
    sufficientemente grande per poter contenere, nel caso peggiore, i
    dati di tutti gli utenti salvati in \texttt{usersFile}. Il programma dovrà poi
    chiedere all’utente il minimo punteggio e, utilizzando la funzione
    \texttt{readPoints()}, leggere da file e memorizzare in V i dati degli
    utenti che hanno almeno il punteggio minimo specificato. Il
    programma infine deve stampare a video il nome ed il punteggio
    degli utenti contenuti in V se e solo se il nome comincia per “Me”.

\end{Exercise}

\begin{Answer}[ref={ex:telefoni}]
  	\lstinputlisting[style=cstylesmall]{src/telefoni.c}
\end{Answer}